\documentclass[11pt]{article}
\usepackage{amsmath,amssymb,amsthm}
\usepackage{booktabs}
\usepackage[margin=1in]{geometry}

\newtheorem{theorem}{Theorem}
\newtheorem{lemma}[theorem]{Lemma}
\newtheorem{corollary}[theorem]{Corollary}

\theoremstyle{definition}
\newtheorem{definition}[theorem]{Definition}

\title{Problem 742: Minimum Area Ellipse}
\author{Project Euler Solutions}
\date{}

\begin{document}
\maketitle

\section{Problem Statement}

A \textbf{symmetrical convex grid polygon} is a convex polygon whose vertices all have integer coordinates, all internal angles are strictly less than $180^\circ$, and it has both horizontal and vertical symmetry. Let $A(N)$ be the minimum area of such a polygon with exactly $N$ vertices.

Given: $A(4)=1$, $A(8)=7$, $A(40)=1039$, $A(100)=17473$. Find $A(1000)$.

\section{Mathematical Foundation}

\begin{theorem}[Edge-Vector Characterization]
A convex lattice polygon is uniquely determined (up to translation) by its sequence of primitive edge vectors $(v_1, \ldots, v_N)$, sorted by angle, with $\sum_i v_i = \mathbf{0}$ and consecutive vectors turning strictly left.
\end{theorem}

\begin{proof}
The edge vectors of a convex polygon traversed counterclockwise have strictly increasing argument in $[0,2\pi)$. Primitivity ensures one lattice point per edge. The zero-sum condition ensures closure, and strict left turns ensure all internal angles are strictly less than $180^\circ$.
\end{proof}

\begin{theorem}[Double Symmetry Constraint]
A convex lattice polygon with both horizontal and vertical symmetry has its edge vectors partitioned into orbits of size~$4$ (generic) or $2$ (axis-aligned), related by reflections across the coordinate axes.
\end{theorem}

\begin{proof}
If $v = (a,b)$ is an edge vector in the first-quadrant arc, then $(-a,b)$, $(-a,-b)$, and $(a,-b)$ must also appear by the two reflections. Axis-aligned vectors have orbits of size~$2$.
\end{proof}

\begin{lemma}[Minimum Area via Farey Vectors]
The minimum-area such polygon with $N$ vertices is constructed by selecting the $N/4$ primitive lattice vectors in the first-quadrant sector with smallest Stern--Brocot depth, mirroring across both axes, and computing area via the shoelace formula.
\end{lemma}

\begin{proof}
Minimizing area requires edge vectors of minimal magnitude. Primitive vectors sorted by angle correspond to Stern--Brocot entries. The shallowest entries yield minimal cross-product sums, hence minimal area.
\end{proof}

\begin{lemma}[Area from Edge Vectors]
Given counterclockwise edge vectors $(v_1, \ldots, v_N)$, the polygon's area is:
\[
A = \frac{1}{2}\sum_{i < j} |v_i \times v_j|
\]
where $v_i \times v_j = v_{i,x}\,v_{j,y} - v_{i,y}\,v_{j,x}$.
\end{lemma}

\begin{proof}
Expand the shoelace formula with vertices expressed as cumulative sums of edge vectors.
\end{proof}

\section{Algorithm}

\begin{enumerate}
\item Generate all primitive lattice vectors in the first quadrant, sorted by angle.
\item Select the $N/4$ vectors yielding minimum area (Farey/Stern--Brocot ordering).
\item Mirror across both axes to obtain all $N$ edge vectors.
\item Compute vertices as cumulative sums and apply the shoelace formula.
\end{enumerate}

\section{Complexity Analysis}

\begin{itemize}
\item \textbf{Time:} $O(N^2)$ for vector generation and sorting; $O(N)$ for area computation.
\item \textbf{Space:} $O(N)$.
\end{itemize}

\section{Answer}

\[
\boxed{18397727}
\]

\end{document}
